\section{Aufbau der lokalen Kubernetes-Testumgebung}

Nachdem Codetask so angepasst wurde, dass es in einen Kubernetes-Cluster deployt werden kann, musste dies ausprobiert und weiterentwickelt werden, um Features von Kubernetes einbauen zu können.
Zu diesem Zweck sollte eine Kubernetes-Testumgebung erstellt werden, die einen Kubernetes-Cluster lokal auf dem Rechner simuliert.
Ziel dabei war, dass die Studierenden, die an diesem Projekt entwickeln, ihre Veränderungen testen können, ohne dass sie erst lernen müssen, wie ein Kubernetes-Cluster funktioniert bzw. verwendet werden muss.
Dabei war es mir auch wichtig, die Einrichtung so einfach wie möglich zu gestalten, damit auf den Rechnern der Studierenden so wenig wie möglich eingerichtet werden muss und diese schneller mit der Entwicklung beginnen können.
In der Vergangenheit hat das lokale Deployment von Codetask den Anfang sehr ausgebremst.
Ein weiteres Ziel ist es, die Manifeste zu erweitern, um diese in Kubernetes besser nutzbar zu machen.

\subsection{Verwendete Technologien/Werkzeuge}

\subsubsection{kind}

Kind ist ein Tool, mit dem sich ein Kubernetes-Cluster lokal auf dem Rechner über Docker simulieren lässt.
Was Kind dabei anders macht als andere Tools, die dasselbe tun, z. B. das weitverbreitete Minikube, ist, dass Kind nicht selbst in Docker läuft, sondern jeder Node in diesem Cluster einen eigenen Container hat.
Dadurch verbraucht Kind viel weniger RAM, was mit meinem Ziel zusammenpasst, dass der lokale Cluster auf jedem Rechner eines Studierenden nutzbar ist.
Meiner persönlichen Meinung nach ist es auch besser, einen Cluster zu simulieren, bei dem jeder Node auf einem anderen Server existiert, da jeder Node in einem eigenen Container isolierter ist.
\subsubsection{Traefik}

Um ein Ingress zunutzen braucht ein Kubernetes Cluster zunächst ein Ingress Controller der als Loadbalancer für den Ingress verwendet wird.
Ich habe aus zwei Gründen mich für Traefik entschieden. Der erste Grund ist das der nginx Ingress Controller im März 2026 von Kubernetes stillgelegt wurde \parencite{kubernetes_ingress_nginx_retirement_2025}.
Der andere Grund ist das Traefik schon in den Compose Datein von der Staging und Production Umgebung als Reverse Proxy genutzt wurde.
\subsection{Umsetzung / Konfiguration}

Kind kann einfach mit einem Install-Tool wie Chocolate heruntergeladen werden.
Um Kind mit mehreren Nodes laufen zu lassen, wurde eine Kind-Konfiguration erstellt, die beim Starten eines Kind-Clusters mitgegeben wird.
Der lokale Cluster hat einen Master-Node und zwei Worker-Nodes. 
Diese Aufteilung dient nur dazu, die Nutzung mehrerer Nodes zu simulieren und zu verhindern, dass die gesamte Last in einem Docker-Container liegt. 

Sobald der Cluster hochgefahren ist, müssen die Helm-Charts in den Cluster geladen werden. 
Danach werden die Images der eigenen Services in ein In-Cluster-Registry geladen. Anschließend wird ein Rollout der Deployments ausgeführt, damit die Pods die neuen Container wahrnehmen.

Damit die Studierenden den Cluster nicht selbst auf ihren lokalen Geräten einrichten müssen, wurde ein Automatisierungsskript entworfen. Für Windows ist dieses Skript in PowerShell geschrieben, während es für macOS und Linux in Shell geschrieben ist.
Das Automation-Script richtet den Kind-Cluster auf dem lokalen Gerät ein, lädt die Helm-Charts in den Cluster und erstellt daraufhin die Images, die in das In-Cluster-Registry gepusht werden. Zum Schluss führt es ein Rollout aus, damit die Deployments die Pods mit den Containern neu starten. 

Durch den lokalen Cluster konnten die eigenen Helm-Charts zunächst ausprobiert und angepasst werden.
Da im lokalen Cluster nur die Helm-Charts und der Code ausprobiert werden müssen und der Cluster nur eine kurze lebig ist, gibt es von jedem Service nur einen Pod.
Das hat auch weiterhin den Nutzen, dass auf den lokalen Geräten der Studierenden, die sich mit diesem Projekt befassen, weniger Ressourcen benötigt werden.

Die erste Anpassung, die der lokale Cluster an den Helm-Charts vorgenommen hat, war das Hinzufügen des Ressourcenmanagements.
In Kubernetes kann für die Pods angegeben werden, wie viel Ressourcen mindestens reserviert werden können und wie viel ein Pod maximal haben darf. 
Zu diesen Ressourcen gehören die benötigte RAM- und CPU-Kapazität. Dadurch kann Kubernetes die Pods besser verwalten, da nun bekannt ist, wie viel Ressourcen ein Pod mindestens braucht, es aber auch ein Limit gibt, damit kein Pod die Ressourcen anderer abziehen kann.
Um zu entscheiden, wie viel RAM und CPU jeder Pod bekommt, wurde ein Monitoring eingesetzt.

Die weitere Anpassung, die vorgenommen wurde, ist, dass Probes für jeden Helm-Chart hinzugefügt werden.
Dabei wurde für jeden Helm-Chart eine StartupProbe und eine ReadinessProbe eingerichtet. 
Eine StartupProbe prüft, ob der Container im Pod hochgefahren ist und die Anwendung gestartet wurde. Solange die StartupProbe nicht bestanden wurde, werden die anderen Probes blockiert.
Wenn eine StartupProbe zu oft fehlschlägt, wird der gesamte Pod neu gestartet.
Eine ReadinessProbe prüft, ob der Container Requests bearbeiten kann, d. h., ob die Anwendung selbst läuft und externe Services noch erreichbar sind, ohne die die komplette Anwendung nicht funktioniert.
Wird eine ReadinessProbe nicht bestanden, wird der Pod auf „NotReady” gesetzt. Dadurch wird der Traffic auf die anderen Pods umgeleitet, bis die ReadinessProbe wieder bestanden wird oder der Pod neu gestartet wird.

Dabei schickt Kubernetes die Probes als Request an die Pods selbst.
Beim Frontend muss nur dem Webserver ein einfacher Request geschickt werden, der mit einem 200-Statuscode beantwortet werden muss.
Bei der Haupt-API, der Check-API und der AI-API mussten eigene Healthcheck-Endpunkte implementiert werden. Diese antworten mit einem 200, wenn sie erreicht werden.
Die MongoDB wird über TCP-Sockets angefragt.

Damit die Service-Frontends, die Haupt-API und die AI-API von außen erreichbar sind, musste für jedes dieser Helm-Charts ein Ingress erstellt werden. 
Da geplant ist, Traefik als Ingress-Controller zu nutzen, wird dessen Helm-Chart auch schon in den lokalen Cluster geladen, um näher an der Produktionsumgebung zu bleiben.
Da die Nodes in Docker-Containern sind, musste über die kind config ein externPortMapping eingerichtet werden, damit die Ingress-Ports von außen über localhost erreichbar sind.
\subsection{Probleme und Anpassungen}
Beim Einrichten des lokalen Clusters sind Probleme aufgetreten.
Das erste Problem das entstanden ist war das lokale Registry der Images. Zwar hat kind eine eigene Lösung dafür, diese funktoniert aber unter Windows nur über den WSL, was gegen das Ziel sprach, dass der Lokale Cluster schnell und simple eingerichtet wird sprach, 
zum weitern arbeitet kind aktuell an einer bessern build-in Lösung für das lokale Registry \parencite{kind_local_registry}.
Daher wurde zunächst das Feature benutzt, bei dem die Images einfach in Kind geladen werden können. Dies war jedoch aufgrund der drei Nodes speicher- und RAM-intensiv, da die Images jeweils in alle drei Container geladen wurden.
Schlussendlich wurde sich daher für ein In-Cluster-Registry entschieden, wodurch das Container Registry in der lokalen Umgebung im Cluster ist.

Ein weiteres Problem entstand durch die Probes. Da die Healthchecks von den Pods selbst durchgeführt werden und die Pods im Kubernetes-Netzwerk eine eigene IP-Adresse haben, standen sie nicht in der Whiteliste der Play-Config, wodurch die Probes ständig fehlschlugen.
Da die Pods kurzlebig sind und bei einem Neustart ihre IP-Adresse verändern, konnte diese auch nicht in die Play-Config hinzugefügt werden. Daher musste in die White List der Play Config eine Wildcard eingefügt werden, um alle Pods zu erlauben. Da die Pods nur intern im Kubernetes-Netzwerk erreichbar sind und nur von außen über Ingress, entsteht dadurch aber kein Sicherheitsrisiko.

Bei der Einrichtung der Ingress gab es für den lokalen Cluster auch Probleme durch den Hostnamen. Da alle drei Services zunächst als Host „localhost” hatten, mussten diese mit Prefixen, die abgeschnitten wurden, unterschieden werden.
Die Präfixe wurden abgeschnitten, was zu Problemen bei den WebSockets der ai-API geführt hat. Dadurch wurden unterschiedliche Hosts benötigt.
Damit dennoch localhost als Basis verwendet werden kann und die Studierenden nichts in ihren Host-Dateien verändern müssen, startet die Automation zusätzlich einen DNS-Resolver mit einer DNS-Wildcard, damit verschiedene Hosts verwendet werden können.
Damit es plattformübergreifend ist, läuft dieser DNS-Resolver in Docker.
\subsection{Ergebnisse der Umsetzung}
Durch die Umsetzung haben die Studierenden einen lokalen Cluster, in dem sie ihre Implementierungen in einer produktionsechten Umgebung ausprobieren können.
Die Automation der Cluster-Einrichtung bringt nicht nur den Vorteil, dass die Studierenden den Cluster leicht auf ihren Geräten einrichten können, ohne Kubernetes-Kenntnisse zu benötigen, sondern der Cluster muss dadurch auch nur temporär auf diesen Geräten vorhanden sein, wenn er gerade benötigt wird. 
Dadurch können die Studierenden beim Entwickeln den Cluster mit nur einem Befehl starten, ihre Implementierungen testen und nach Abschluss der Arbeit diesen wieder mit einem Cleanup-Befehl entfernen, bis sie den lokalen Cluster wieder benötigen.

Durch die Ressourcenverwaltung können nun Pods besser auf den Nodes verteilt werden, zusätzlich wird auch verhindert das ein Service alle Ressourcen für sich selbst beanspruchen kann.
Zusammen mit den Probes und den RollingUpdates erzielt das Codetask bei einem Update eine Zero Downtime.

Durch den Ingress sind nun drei Services nach außen erreichbar mit http, während die anderen Services nur im Cluster selbst erreichbar bleiben.